%% kapa: exam=1 round=29 year=2018 period=1 form=A subject=civil_law subject_ko=민법 seq=34 number=34
%% source: Q-Net 감정평가사 기출문제 (cst00309) artlSeq=5202808 file=extracted/1008/5202808/2018년 제29회 감정평가사 1차시험 기출문제/2018년 제29회 감정평가사 1차시험 기출문제/1교시 A형.pdf
%% conversion: kordoc 4.24.51 -> tools/gen1.py (auto); review=proofread-sonnet; flags=none
\begin{question}{34}
토지소유권에 관한 설명으로 옳지 않은 것은? (다툼이 있으면 판례에 따름)
\begin{choices}
  \item 토지의 소유권은 정당한 이익있는 범위 내에서 토지의 상하에 미친다.
  \item 명인방법을 갖춘 수목의 집단은 토지의 구성부분이 아니다.
  \item 토지가 해면 아래에 잠김으로써 포락될 당시를 기준으로 원상복구가 불가능한 상태에 이르면 종전의 소유권은 영구히 소멸된다.
  \item 토지등기부의 표제부에 토지의 면적이 실제와 다르게 등재되어 있으면, 이러한 등기는 해당 토지를 표상하는 등기로서 효력이 없다.
  \item 토지 1필지의 공간적 범위를 특정하는 것은 지적도나 임야도의 경계이지 등기부의 표제부나 임야대장ㆍ토지대장에 등재된 면적이 아니다.
\end{choices}
\end{question}
\answer{4}
